\documentclass[12pt]{article}
\usepackage[T1]{fontenc}
\usepackage[utf8]{inputenc}
\usepackage{lmodern}
\usepackage{textcomp}
\usepackage{enumitem}
\usepackage{geometry}
\geometry{margin=1in}
\begin{document}
\begin{center}
Answer Key - History - Graduate Level Assessment 25 \
\small (Answers are ordered exactly as in the question paper.)
\end{center}

\section*{Multiple Choice}
\begin{enumerate}[label=\arabic*.]
\item \textbf{Answer:} A
\\ \textit{Options:} A) North America; B) Australia; C) Central America; D) Arctic Region; E) Middle East
\\ \textit{Note:} The Hoabinhian tradition, characterized by the manufacture of stone tools from chipped pebbles, is primarily associated with prehistoric cultures in Southeast Asia rather than North America, making the provided answer incorrect; however, if the focus were on adaptations of similar stone tool technologies in varied regions, one could argue for a broader interpretation of tool-making practices.
\item \textbf{Answer:} E
\\ \textit{Options:} A) Hohokam; B) Folsom; C) Adena; D) Ancestral Puebloan; E) Mississippian
\\ \textit{Note:} The Mississippian culture, known for their construction of large temple mounds, developed advanced agricultural practices that centered around the cultivation of maize and squash, which were staple crops in their diet and economy.
\item \textbf{Answer:} A
\\ \textit{Options:} A) office; landfill; B) dining area; garbage disposal; C) bathroom; incinerator; D) kitchen; garden; E) playground; waste treatment plant
\\ \textit{Note:} An office contains primary refuse, such as documents and materials directly related to its operations, while a landfill is designed to manage secondary refuse, which includes waste that has been discarded or processed from various sources.
\item \textbf{Answer:} A
\\ \textit{Options:} A) the study of plant and animal life from a particular period in history.; B) a study of the spatial relationships among artifacts, ecofacts, and features.; C) the study of human behavior and culture.; D) a technique to determine the age of artifacts based on their position in the soil.; E) a dating technique based on the decay of a radioactive isotope of bone.
\\ \textit{Note:} Taphonomy is the branch of paleontology that examines the processes of decay, preservation, and the fossilization of organisms, thereby providing insights into the plant and animal life from specific historical periods.
\item \textbf{Answer:} A
\\ \textit{Options:} A) Mogollon; B) Navajo; C) Clovis; D) Paleo-Indians; E) Hopewell
\\ \textit{Note:} The Mogollon culture is recognized for its development of permanent settlements and agricultural practices in the Southwest, distinguishing them as a primarily sedentary society compared to their nomadic counterparts.
\end{enumerate}

\section*{True / False}
\begin{enumerate}[label=\arabic*.]
\setcounter{enumi}{5}
\item \textbf{Answer:} True
\\ \textit{Options:} A) True; B) False
\\ \textit{Note:} According to the passage: irrational and backward
\item \textbf{Answer:} True
\\ \textit{Options:} A) True; B) False
\\ \textit{Note:} According to the passage: to stay, so long as there was at least an "indirect quid pro quo" for the work he did
\item \textbf{Answer:} True
\\ \textit{Options:} A) True; B) False
\\ \textit{Note:} According to the passage: 20 th century
\item \textbf{Answer:} False
\\ \textit{Options:} A) True; B) False
\\ \textit{Note:} The correct answer is: Janissary, Sipahi, Akıncı and Mehterân
\item \textbf{Answer:} False
\\ \textit{Options:} A) True; B) False
\\ \textit{Note:} The correct answer is: Georgics
\end{enumerate}

\section*{Long Form Response}
\begin{enumerate}[label=\arabic*.]
\setcounter{enumi}{10}
\item \textbf{Answer:} seventy translators
\\ \textit{Note:} Expected answer: seventy translators
\item \textbf{Answer:} Hellenistic
\\ \textit{Note:} Expected answer: Hellenistic
\end{enumerate}

\vfill
\noindent\textit{End of Answer Key}
\end{document}